\begin{table}[htbp]\centering\footnotesize
\caption{La vara del artículo 61 de la Ley de Vivienda contra el Ramo 15 (mdp)}\label{cua:vivienda}
\begin{tabular}{lrrrr}
\toprule
Necesidad & Hogares & UMA & Costo unitario (pesos) & Monto (mdp) \\
\midrule
Mejora de vivienda & 849.226 & 25 & 89.155,5 & 75.713,2 \\
Ampliación de vivienda & 771.091 & 50 & 178.311,0 & 137.494,0 \\
Vivienda nueva & 36.246 & 100 & 356.622,0 & 12.926,1 \\
\midrule \textbf{Requerimiento declarado} & & & & \textbf{226.133,3} \\
Ramo 15 completo, proyecto 2027 & & & & 37.873,6 \\
\textbf{El ramo como proporción del requerimiento} & & & & \textbf{16,7\,\%} \\
\bottomrule
\end{tabular}
\\[2pt]\begin{minipage}{0.92\textwidth}\scriptsize Fuente: Gaceta Parlamentaria 7121, anexo N, estimación de la Secretaría de Bienestar; y \texttt{datos/ramos.csv} para el Ramo 15. \textbf{Años de pesos distintos}: los costos unitarios provienen de las Reglas de Operación y de la UMA de 2026, y el presupuesto está en pesos de 2027. \textbf{El Ramo 15 no es el perímetro del gasto en subsidios de vivienda}: es el ramo completo, y CONAVI es un organismo descentralizado. La comparación es de orden de magnitud.\end{minipage}
\end{table}
